% =============================================================================
%  Notas de clase — Sesión 10: Revisión del Examen Parcial 01
%  Programación Avanzada — Universidad Panamericana
%  Compilar: pdflatex sesion10_revision_examen.tex
% =============================================================================
\documentclass[12pt, oneside]{article}
\usepackage{progavanzada}

\begin{document}

\portada{Sesión 10}{Revisión del Examen Parcial 01}{2026}

\tableofcontents
\newpage

% =============================================================================
\section{¿De qué trata esta sesión?}
% =============================================================================

Esta sesión \textbf{no introduce material nuevo}.
El objetivo es revisar juntos el examen que resolviste en la sesión anterior:
ver qué salió bien, qué salió mal, y \textbf{por qué}.

El profesor regresará los exámenes calificados y revisará cada pregunta
en el pizarrón.  Tu tarea durante la sesión es escuchar, comparar con tu
propio examen y anotar lo que no quedó claro.

\begin{recuerda}
  Revisar un examen no es solo saber cuántos puntos obtuviste.
  Es la oportunidad de cerrar los huecos conceptuales antes de que el curso
  avance y esos huecos se vuelvan más costosos.
\end{recuerda}

% =============================================================================
\section{¿Para qué sirve esta pausa?}
% =============================================================================

Aprender a programar no es lineal: a veces crees que entendiste algo,
lo aplicas en el examen y descubres que había un detalle que se te escapó.
Ese momento de \textbf{disonancia} —cuando lo que pensabas que sabías choca
con el resultado— es uno de los más valiosos del proceso de aprendizaje.

Esta sesión existe precisamente para aprovechar ese momento.

\begin{consejo}
  Cuando el profesor revise una pregunta que tú contestaste diferente,
  no te limites a copiar la respuesta correcta.
  \textbf{Pregunta por qué}: ¿qué suposición estaba mal?
  ¿Qué parte del proceso fallaste?
  Anotar eso es mucho más útil que memorizar la respuesta.
\end{consejo}

% =============================================================================
\section{Espacio para tus notas de revisión}
% =============================================================================

Usa este espacio para anotar los errores que encontraste en tu examen y
la explicación correcta.

\vspace{0.4cm}
\noindent\textbf{Pregunta / tema:} \hrulefill

\vspace{1.8cm}
\noindent\textbf{Mi error:} \hrulefill

\vspace{1.8cm}
\noindent\textbf{Lo correcto y por qué:}

\vspace{3cm}
\noindent\rule{\textwidth}{0.4pt}

\vspace{0.6cm}
\noindent\textbf{Pregunta / tema:} \hrulefill

\vspace{1.8cm}
\noindent\textbf{Mi error:} \hrulefill

\vspace{1.8cm}
\noindent\textbf{Lo correcto y por qué:}

\vspace{3cm}
\noindent\rule{\textwidth}{0.4pt}

\vspace{0.6cm}
\noindent\textbf{Pregunta / tema:} \hrulefill

\vspace{1.8cm}
\noindent\textbf{Mi error:} \hrulefill

\vspace{1.8cm}
\noindent\textbf{Lo correcto y por qué:}

\vspace{3cm}
\noindent\rule{\textwidth}{0.4pt}

\end{document}
